%=======================================================================
%
% ------------------------------------------------------------------------
% ------------------------------------------------------------------------
% abnTeX2: Modelo de Trabalho Academico (tese de doutorado, dissertacao de
% mestrado e trabalhos monograficos em geral) em conformidade com 
% ABNT NBR 14724:2011: Informacao e documentacao - Trabalhos academicos - Apresentacao
% ------------------------------------------------------------------------
% ------------------------------------------------------------------------
% Autor:  Jonas Alessi(alessi.jonas@gmail.com
% Versão: 10 de Junlho 2013.
% Edição: TexStudio
% LaTeX:  abnTeX2
%
%=======================================================================

\documentclass[
	% -- opções da classe memoir --
	12pt,				% tamanho da fonte	
	oneside,          % não imprimir em verso e anverso, oposto do twoside 
	a4paper,			% tamanho do papel. 
	% -- opções da classe abntex2 --
	chapter=TITLE,		% títulos de capítulos convertidos em letras maiúsculas
	section=TITLE,		% títulos de seções convertidos em letras maiúsculas
	subsection=TITLE,	% títulos de subseções convertidos em letras maiúsculas
	%subsubsection=TITLE,% títulos de subsubseções convertidos em letras maiúsculas
	% -- opções do pacote babel --
	english,			% idioma adicional para hifenização
	brazil,				% o último idioma é o principal do documento
	]{abntUnoesc}


% ---
% Pacotes fundamentais 
% ---
\usepackage{cmap}				% Mapear caracteres especiais no PDF
\usepackage{times}			    % Usa a fonte Latin Modern			
\usepackage[T1]{fontenc}		% Selecao de codigos de fonte.
\usepackage[utf8]{inputenc}		% Codificacao do documento (conversão automática dos acentos)
\usepackage{lastpage}			% Usado pela Ficha catalográfica
\usepackage{indentfirst}		% Indenta o primeiro parágrafo de cada seção.
\usepackage{color}				% Controle das cores
\usepackage{graphicx}			% Inclusão de gráficos

% ---
% Pacotes de citações
% ---
\usepackage[brazilian,hyperpageref]{backref}	 % Paginas com as citações na bibl
\usepackage[alf]{abntex2cite}					 % Citações padrão ABNT
\usepackage{lipsum}							   	 % para geração de dummy text

% Configurações de aparência do PDF final
% alterando o aspecto da cor azul
\definecolor{blue}{RGB}{41,5,195}
% --- 
% Espaçamentos entre linhas e parágrafos 
% --- 
% O tamanho do parágrafo é dado por:
\setlength{\parindent}{1.25cm}
\linespread{1.5}

% Controle do espaçamento entre um parágrafo e outro:
\setlength{\parskip}{0.2cm}  % tente também \onelineskip
%Espaçamento padrão entrelinhas 
\renewcommand{\backref}{}

% ---
% compila o indice
% ---
\makeindex
% ---

%----Include da capa é fora do documento 
% ---
% Informações de dados para CAPA e FOLHA DE ROSTO
% ---
\titulo{TITULO:Subtitulo(se houver)}
\autor{NOME DO AUTOR DO TRABALHO}
\local{Joaçaba - SC}
\data{2013}
\orientador{Orientador 1}
\coorientador{Coorientador 2}
%\coorientador{NOME DO COORIENTADOR}
\instituicao{NOME DA INSTITUIÇÃO}
\tipotrabalho{Tese (Doutorado)}
% O preambulo deve conter o tipo do trabalho, o objetivo, 
% o nome da instituição e a área de concentração 
\preambulo{Trabalho de Conclusão de Curso apresentado ao Curso de ..., Área
da Ciências ... da Universidade ... como requisitos parcial à obtenção do grau de ... em ...}
% ---

%---

\begin{document}
%\OnehalfSpacing
% Retira espaço extra obsoleto entre as frases.
\frenchspacing 

% ----------------------------------------------------------
% ELEMENTOS PRÉ-TEXTUAIS
% ----------------------------------------------------------

%--- CAPA ----
\imprimircapa

% --- FOLHA DE ROSTO
\imprimirfolhaderosto
% \imprimirfolhaderosto* (o * indica que haverá a ficha bibliográfica)

% Inserir FOLHA DE APROVAÇÃO
\begin{folhadeaprovacao}

  \begin{center}
    {\ABNTEXchapterfont\large\imprimirautor}

    \vspace*{\fill}\vspace*{\fill}
    {\ABNTEXchapterfont\Large\imprimirtitulo}
    \vspace*{\fill}
    
    \hspace{.45\textwidth}
    \begin{minipage}{.5\textwidth}
        \imprimirpreambulo
    \end{minipage}%
    \vspace*{\fill}
   \end{center}
    
  \begin{flushleft}
  	 Aprovado em 24 de novembro de 2012
  \end{flushleft}
  
  \begin{center}
  BANCA EXAMINADORA
  \end{center}

   \assinatura{\imprimirorientador \\ Orientador} 
   \assinatura{Professor \\ Convidado 1}
   \assinatura{Professor \\ Convidado 2}
   %\assinatura{\textbf{Professor} \\ Convidado 3}
   %\assinatura{\textbf{Professor} \\ Convidado 4}
      
   \begin{center}
    \vspace*{0.5cm}
    {\large\imprimirlocal}
    \par
    {\large\imprimirdata}
    \vspace*{1cm}
  \end{center}
  
\end{folhadeaprovacao}

% DEDICATÓRIA
\begin{dedicatoria}
 \vspace*{\fill}
 \noindent
  \raggedleft
 \begin{minipage}{.54\textwidth}
    Dedico este trabalho a meus pais, fonte de meus conhecimentos  e saber. Graças a eles, tornei-me uma pessoa capaz de lutar para que meus sonhos e objetivos fossem sempre alcançados, sem jamais desanimar. Considero-me forte porque eles me ensinaram a ser forte.
   \end{minipage}
\end{dedicatoria}


% AGRADECIMENTO
\begin{agradecimentos}
aaaaa bbbbb ccccc aaaaaaaaaaaa aa dd  aaa aaa aaaa aaaaaaaaaaa aassss aaaa 
aaaaa bbbbb ccccc aaaaaaaaaaaa aa dd  aaa aaa aaaa aaaaaaaaaaa aassss aaaa
aaaaa bbbbb ccccc aaaaaaaaaaaa aa dd  aaa aaa aaaa aaaaaaaaaaa aassss aaaa
aaaaa bbbbb ccccc aaaaaaaaaaaa aa dd  aaa aaa aaaa aaaaaaaaaaa aassss aaaa
aaaaa bbbbb ccccc aaaaaaaaaaaa aa dd  aaa aaa aaaa aaaaaaaaaaa aassss aaaa
aaaaa bbbbb ccccc aaaaaaaaaaaa aa dd  aaa aaa aaaa aaaaaaaaaaa aassss aaaa
aaaaa bbbbb ccccc aaaaaaaaaaaa aa dd  aaa aaa aaaa aaaaaaaaaaa aassss aaaa
aaaaa bbbbb ccccc aaaaaaaaaaaa aa dd  aaa aaa aaaa aaaaaaaaaaa aassss aaaa
aaaaa bbbbb ccccc aaaaaaaaaaaa aa dd  aaa aaa aaaa aaaaaaaaaaa aassss aaaa
aaaaa bbbbb ccccc aaaaaaaaaaaa aa dd  aaa aaa aaaa aaaaaaaaaaa aassss aaaa
aaaaa bbbbb ccccc aaaaaaaaaaaa aa dd  aaa aaa aaaa aaaaaaaaaaa aassss aaaa
aaaaa bbbbb ccccc aaaaaaaaaaaa aa dd  aaa aaa aaaa aaaaaaaaaaa aassss aaaa
aaaaa bbbbb ccccc aaaaaaaaaaaa aa dd  aaa aaa aaaa aaaaaaaaaaa aassss aaaa
aaaaa bbbbb ccccc aaaaaaaaaaaa aa dd  aaa aaa aaaa aaaaaaaaaaa aassss aaaa
\end{agradecimentos}

%EPÍGRAFE
\begin{epigrafe}
    \vspace*{\fill}
	\begin{flushright}
		\textit{``Não vos amoldeis às estruturas deste mundo, \\
		mas transformai-vos pela renovação da mente, \\
		a fim de distinguir qual é a vontade de Deus: \\
		o que é bom, o que Lhe é agradável, o que é perfeito.\\
		(Bíblia Sagrada, Romanos 12, 2)}
	\end{flushright}
\end{epigrafe}

% RESUMO

% resumo em português
\begin{resumo}
 Segundo a citeonline[3.1-3.2]{NBR6028:2003}, o resumo deve ressaltar o
 objetivo, o método, os resultados e as conclusões do documento. A ordem e a extensão
 destes itens dependem do tipo de resumo (informativo ou indicativo) e do
 tratamento que cada item recebe no documento original. O resumo deve ser
 precedido da referência do documento, com exceção do resumo inserido no
 próprio documento. (\ldots) As palavras-chave devem figurar logo abaixo do
 resumo, antecedidas da expressão Palavras-chave:, separadas entre si por
 ponto e finalizadas também por ponto.

 \vspace{\onelineskip}
    
 \noindent
 \textbf{Palavras-chaves}: latex. abntex. editoração de texto.
\end{resumo}

% resumo em inglês
\begin{resumo}[Abstract]
 \begin{otherlanguage*}{english}
   This is the english abstract.

   \vspace{\onelineskip}
 
   \noindent 
   \textbf{Key-words}: latex. abntex. text editoration.
 \end{otherlanguage*}
\end{resumo}

% lista de ilustrações
\pdfbookmark[0]{\listfigurename}{lof}
\listoffigures*
\cleardoublepage

% lista de quadros
\pdfbookmark[0]{\listQuadroname}{lof}
\listoffigures*
\cleardoublepage

% LISTA DE TABELAS
\pdfbookmark[0]{\listtablename}{lot}
\listoftables*
\cleardoublepage  %-- força proxima pagina

% ---
% inserir lista de abreviaturas e siglas
% ---
\begin{siglas}
  \item[ABNT] Associação Brasileira de Normas Técnicas
  \item[Bras.] Brasileiro
  \item[CNPq] Conselho Nacional de Desenvolvimento Científico e Tecnológico
  \item[IBGE] Instituto Brasileiro de Geografia e Estatística
\end{siglas}
% ---

% ---
% inserir lista de símbolos
% ---
\begin{simbolos}
  \item[$ \Gamma $] Letra grega Gama
  \item[$ \Lambda $] Lambda
  \item[$ \zeta $] Letra grega minúscula zeta
  \item[$ \in $] Pertence
\end{simbolos}
% ---

% SUMARIO
\pdfbookmark[0]{\contentsname}{toc}
\tableofcontents*
\cleardoublepage

% ------------------------------------------------------
% ELEMENTOS TEXTUAIS
% ------------------------------------------------------
\textual
% INTRODUÇÃO
% ----------------------------------------------------------
% Introdução
%Ex: \chapter{TÍTULO A SER IMPRESSO NO CORPO DO TEXTO}{Título no cabeçalho}{Título no Sumario}
% ----------------------------------------------------------
\chapter[INTRODUÇÃO]{INTRODUÇÃO} % se usar o \chapter* ele não vai colocar no sumario
%\addcontentsline{toc}{chapter}{Introdução} inclui manualmente no sumario

------ ------- -------- --  -------- ----- -- -- ----- -- -- --  -- - --- - -- - - - - - -
------ ------- -------- --  -------- ----- -- -- ----- -- -- --  -- - --- - -- - - - - - -
------ ------- -------- --  -------- ----- -- -- ----- -- -- --  -- - --- - -- - - - - - -
------ ------- -------- --  -------- ----- -- -- ----- -- -- --  -- - --- - -- - - - - - -
------ ------- -------- --  -------- ----- -- -- ----- -- -- --  -- - --- - -- - - - - - -
------ ------- -------- --  -------- ----- -- -- ----- -- -- --  -- - --- - -- - - - - - -
------ ------- -------- --  -------- ----- -- -- ----- -- -- --  -- - --- - -- - - - - - -
------ ------- -------- --  -------- ----- -- -- ----- -- -- --  -- - --- - -- - - - - - -
------ ------- -------- --  -------- ----- -- -- ----- -- -- --  -- - --- - -- - - - - - -
------ ------- -------- --  -------- ----- -- -- ----- -- -- --  -- - --- - -- - - - - - -
------ ------- -------- --  -------- ----- -- -- ----- -- -- --  -- - --- - -- - - - - - -
------ ------- -------- --  -------- ----- -- -- ----- -- -- --  -- - --- - -- - - - - - -
------ ------- -------- --  -------- ----- -- -- ----- -- -- --  -- - --- - -- - - - - - -
------ ------- -------- --  -------- ----- -- -- ----- -- -- --  -- - --- - -- - - - - - -
------ ------- -------- --  -------- ----- -- -- ----- -- -- --  -- - --- - -- - - - - - -
------ ------- -------- --  -------- ----- -- -- ----- -- -- --  -- - --- - -- - - - - - -
------ ------- -------- --  -------- ----- -- -- ----- -- -- --  -- - --- - -- - - - - - -
------ ------- -------- --  -------- ----- -- -- ----- -- -- --  -- - --- - -- - - - - - ------- ------- -------- --  -------- ----- -- -- ----- -- -- --  -- - --- - -- - - - - - -
------ ------- -------- --  -------- ----- -- -- ----- -- -- --  -- - --- - -- - - - - - -
------ ------- -------- --  -------- ----- -- -- ----- -- -- --  -- - --- - -- - - - - - -
------ ------- -------- --  -------- ----- -- -- ----- -- -- --  -- - --- - -- - - - - - -
------ ------- -------- --  -------- ----- -- -- ----- -- -- --  -- - --- - -- - - - - - -
------ ------- -------- --  -------- ----- -- -- ----- -- -- --  -- - --- - -- - - - - - -

%FUNDAMENTAÇÃO TEÓRICA
\chapter{FUNDAMENTAÇÃO TEÓRICA}

\textcolor{red}{Apresentar estudos que contemple a temática abordada. Respeitar a autoria,
nas citações diretas e indiretas. Evitar parágrafos muito longos. Evitar seções e
subseções muito curtas.}


%DESENVOLVIMENTO
\chapter{DESENVOLVIMENTO}

%%gerar conteudo para teste(5 paragrafos)

Nam dui ligula, fringilla a, euismod sodales, sollicitudin vel, wisi. Morbi auctor lorem non justo. Nam lacus libero, pretium at, lobortis vitae, ultricies et, tellus. Donec aliquet, tortor sed accumsan bibendum, erant montes, nascetur ridiculus mus. Aliquam tincidunt urna. Nulla ullamcorper vestibulum turpis. Pellentesque cursus luctus mauris. \cite{LUCKMANN2008}

Nam dui ligulna, vitae ornare odio metus a mi. Morbi ac orci et nisl hendrerit mollis. Suspendisse ut massa. Cras nec ante. Pellentesque a nulla. Cum sociis natoque penatibus et magnis dis parturient montes, nascetur ridiculus mus. Aliquam tincidunt urna. Nulla ullamcorper vestibulum turpis. Pellentesque cursus luctus mauris. \cite{Telles1984}

Nao metus a mi. Morbi ac orci et nisl hendrerit mollis. Suspendisse ut massa. Cras nec ante. Pellentesque a nulla. Cum sociis natoque penatibus et magnis dis parturient montes, nascetur ridiculus mus. Aliquam tincidunt urna. Nulla ullamcorper vestibulum turpis. Pellentesque cursus luctus mauris. \cite{abntex2-wiki-como-customizar}

Nam duina, vitae ornare odio metus a mi. Morbi ac orci et nisl hendrerit mollis. Suspendisse ut massa. Cras nec ante. Pellentesque a nulla. Cum sociis natoque penatibus et magnis dis parturient montes, nascetur ridiculus mus. Aliquam tincidunt urna. Nulla ullamcorper vestibulum turpis. Pellentesque cursus luctus mauris. \cite{abntex2modelo}

Nam dui ligulaquet magna, vitae ornare odio metus a mi. Morbi ac orci et nisl hendrerit mollis. Suspendisse ut massa. Cras nec ante. Pellentesque a nulla. Cum sociis natoque penatibus et magnis dis parturient montes, nascetur ridiculus mus. Aliquam tincidunt urna. Nulla ullamcorper vestibulum turpis. Pellentesque cursus luctus mauris. \cite{memoir}

\section{SEÇÃO}

\lipsum[1-1]
%\cite{abntex2-wiki-como-customizar} 
\subsection{Subseção de desenvolvimento}

\lipsum[1-1]
%\cite{abntex2modelo
\subsection{Tabela}
%%% EXEMPLO TABELA 
\begin{table}[htb]
	
	\label{tab-inicial}
\begin{center}
	\changecaptionwidth 
	\captionwidth{12.2cm}
    \caption[Título na lista de tabelas]{Título da tabela}
    \begin{tabular}{p{2.6cm}lp{6.0cm}lp{2.25cm}lp{3.40cm}}
  		 \hline
    	\textbf{Título} & \textbf{Título} & \textbf{Título} & \textbf{Título} \\ \hline
   		 tabela & tabela & tabela & tabela\\
         tabela & tabela & tabela & tabela\\
  		 tabela & tabela & tabela & tabela\\ \hline
    \end{tabular}
  \legend{Fonte: \citeonline{livro-unoesc}}
\end{center}
\end{table}

\lipsum[1-1]
%\cite{Telles1984}
\subsection{Figuras}
\lipsum[1-1]
\begin{figure}[htb]
	\begin{center}
		\changecaptionwidth 
		\captionwidth{5.2cm}
		\caption{\label{fig_grafico}Small Tux}
	    \includegraphics[scale=0.2]{figuras/tux.png}
	    \legend{Fonte:\citeonline[p. 44]{livro-unoesc}}
	\end{center}
\end{figure}

\subsubsection{Seção quaternária}
Nam dui ligula, fringilla a, euismod sodales, sollicitudin vel, wisi. Morbi auctor lorem non justo dfsdf sdfsdfsd sdfsdfsdf.

Segundo  \citeonline[p. 94]{livro-unoesc},
\begin{citacao}
Compreende trecho transcrito que apresenta mais de três linhas; mantém-se o discurso do texto original; destaca-se em blocos, espaço simples,com recuo de 4 cm a partir da margem esquerda, com letra menor que a do texto original;sugere-se usar tamanho 10.
\end{citacao}

\lipsum[1-1]
%\cite{LUCKMANN2008}

\section {Listas}

Lorem ipsum dolor sit amet, consectetuer adipiscing elit. Ut purus elit, vestibulum ut, placerat ac, adipiscing vitae, felis. Curabitur dictum gravida mauris. Nam arcu libero, nonummy eget, consectetuer id, vulputate a, magna.

\begin{itemize}
\item Nam dui ligula, fringilla.
\item Nam dui ligula, fringilla.
\item Nam dui ligula, fringilla.
\end{itemize}

Lorem ipsum dolor sit amet, consectetuer adipiscing elit. Ut purus elit, vestibulum ut, placerat ac, adipiscing vitae, felis. Curabitur dictum gravida mauris. Nam arcu libero, nonummy eget, consectetuer id, vulputate a, magna.

\begin{enumerate}
 \item enumA
 \item enumB
\end{enumerate}


%RESULTADOS
\chapter{RESULTADOS}

\lipsum[1-1]



%CONCLUSÃO
\chapter{CONCLUSÃO}

\lipsum[1-1]

% ----------------------------------------------------------
% ELEMENTOS PÓS-TEXTUAIS
% ----------------------------------------------------------
\postextual

% ----------------------------------------------------------
% Referências bibliográficas
% ----------------------------------------------------------

\bibliography{referencias}

%Apêndices
\begin{apendicesenv}

\chapter*{Título do Apêndice A}
\lipsum[1-1]


\end{apendicesenv}

%Anexos
\begin{anexosenv}
\chapter{Título do Anexo A}

\lipsum[1-1]


\end{anexosenv}

%---------------------------------------------------------------------
% INDICE REMISSIVO
%---------------------------------------------------------------------

\printindex
\end{document}